\documentclass[fleqn]{beamer}

\mode<presentation> {
\usetheme{Madrid}
\usecolortheme{dolphin}
}

\usepackage[utf8]{inputenc}
\usepackage[brazilian]{babel}
\usepackage{amsmath}
\usepackage{amssymb}
\usepackage{booktabs}
\usepackage{listings}
\usepackage{xcolor}
\usepackage{tikz}
\usepackage{graphicx}

\lstset{
	basicstyle=\ttfamily\scriptsize,
	breaklines=true,
	keywordstyle=\color{blue!60!black},
	commentstyle=\color{gray},
	stringstyle=\color{green!45!black},
	showstringspaces=false,
	columns=fullflexible,
	frame=single,
	framesep=4pt,
	aboveskip=6pt,
	belowskip=2pt,
}

% ---------- Figuras reutilizáveis ----------

% Paisagem unidimensional de um problema de maximização:
% três máximos locais (x = 1,8; 5,8; 9,0), sendo 5,8 o global.
\newcommand{\paisagem}{%
	\draw[->] (-0.2,0) -- (9.8,0) node[right] {\scriptsize soluções};
	\draw[->] (0,-0.1) -- (0,4.2) node[above] {\scriptsize $f(x)$};
	\draw[thick,blue!60!black] plot[smooth] coordinates {%
		(0.3,1.0) (1.0,2.0) (1.8,2.6) (2.6,2.0) (3.6,0.8) (4.6,1.6)
		(5.8,3.6) (6.8,2.4) (7.6,1.8) (8.4,2.4) (9.0,2.6) (9.6,2.2)};
}

% Linha de um vetor binário (#1 = coordenada y, #2 = lista de bits)
\tikzset{bit/.style={draw,minimum size=0.6cm,inner sep=0pt,font=\small},
	bit0/.style={bit,fill=white}, bit1/.style={bit,fill=green!25}}
\newcommand{\vetor}[2]{%
	\foreach \v [count=\i] in {#2}{%
		\node[bit\v] at (\i*0.65,#1) {\v};
	}%
}
% Contorno em destaque de uma posição do vetor (#1 = y, #2 = índice)
\newcommand{\destaque}[2]{%
	\draw[red,very thick] ({#2*0.65-0.3},{#1-0.3}) rectangle ({#2*0.65+0.3},{#1+0.3});
}

\title[Aula 7 -- Busca Local]{Busca Local e Vizinhança}
\subtitle{Melhorando uma solução construtiva até um ótimo local}
\author[Santi]{Prof. \'Everton Santi}
\institute[UFRN/ECT]{ECT2524 -- Introdução à Otimização \\ Escola de Ciências e Tecnologia -- UFRN}
\date{Aula 7}

\begin{document}

\begin{frame}
	\titlepage
\end{frame}

\begin{frame}
	\frametitle{Sumário}
	\tableofcontents
\end{frame}

%------------------------------------------------------------
\section{Motivação}
%------------------------------------------------------------

\begin{frame}
	\frametitle{Da solução construtiva à solução melhorada}
	\begin{itemize}
		\item Uma heurística construtiva produz uma solução viável com
			rapidez, mas sem garantia de qualidade;
		\item Os GAPs obtidos na tarefa anterior (vizinho mais próximo e
			construção aleatória) mostram que há margem para melhoria;
		\item Ideia central: partir da solução construtiva e aplicar
			pequenas modificações, mantendo apenas as que melhoram o
			custo;
		\item O processo termina em uma solução que nenhuma modificação
			do tipo escolhido consegue melhorar: um \textbf{ótimo local};
		\item Este procedimento é a \textbf{busca local}; as modificações
			permitidas definem a \textbf{vizinhança}.
	\end{itemize}
\end{frame}

\begin{frame}
	\frametitle{Roteiro da aula}
	\begin{itemize}
		\item Espaço de soluções, ótimo local e ótimo global;
		\item Vizinhança e algoritmo de busca local;
		\item Exploração e explotação;
		\item Vizinhanças para o problema da mochila: inversão de bit,
			inserção e troca;
		\item Vizinhança 2-opt para o TSP;
		\item Tarefa: aplicar a busca local 2-opt às soluções das três
			heurísticas construtivas.
	\end{itemize}
\end{frame}

%------------------------------------------------------------
\section{Espaço de soluções e ótimos locais}
%------------------------------------------------------------

\begin{frame}
	\frametitle{Espaço de soluções e função objetivo}
	\begin{itemize}
		\item Conjunto $S$ das soluções viáveis; função objetivo
			$f : S \to \mathbb{R}$;
		\item Visualização: cada ponto do eixo horizontal é uma solução e
			a altura é o valor de $f$ (problema de maximização);
		\item Em problemas combinatórios, $S$ é finito, mas
			enorme: $2^n$ soluções na mochila e $(n-1)!/2$ rotas no TSP.
	\end{itemize}
	\vspace{2mm}
	\begin{center}
	\begin{tikzpicture}[xscale=0.85,yscale=0.65]
		\paisagem
	\end{tikzpicture}
	\end{center}
\end{frame}

\begin{frame}
	\frametitle{Ótimo local e ótimo global}
	\begin{columns}[T]
	\begin{column}{0.48\textwidth}
	\begin{itemize}
		\item \textbf{Ótimo global}: $x^* \in S$ tal que
			$f(x^*) \geq f(x)$ para todo $x \in S$;
		\item \textbf{Ótimo local} em relação a $\mathcal{N}$:
			$\bar{x}\in S$ tal que $f(\bar{x}) \geq f(x)$ para todo
			$x \in \mathcal{N}(\bar{x})$;
		\item Todo ótimo global é ótimo local, qualquer que seja
			$\mathcal{N}$; a recíproca é falsa;
		\item Para minimização, inverte-se a desigualdade.
	\end{itemize}
	\end{column}
	\begin{column}{0.5\textwidth}
	\begin{tikzpicture}[xscale=0.44,yscale=0.6]
		\paisagem
		\fill[red] (1.8,2.6) circle (3pt) node[above=2pt] {\scriptsize local};
		\fill[red] (9.0,2.6) circle (3pt) node[above=2pt] {\scriptsize local};
		\fill[green!50!black] (5.8,3.6) circle (3pt) node[above=2pt] {\scriptsize global};
		\fill[orange] (3.6,0.8) circle (3pt) node[below=2pt] {\scriptsize mín. local};
	\end{tikzpicture}
	\end{column}
	\end{columns}
\end{frame}

%------------------------------------------------------------
\section{Vizinhança e busca local}
%------------------------------------------------------------

\begin{frame}
	\frametitle{Vizinhança}
	\begin{block}{Definição}
		A vizinhança $\mathcal{N}(x)$ de uma solução $x$ é o conjunto de
		soluções obtidas de $x$ por uma pequena alteração de seus
		componentes, definida por um operador de movimento.
	\end{block}
	\begin{itemize}
		\item Depende da representação da solução e do operador escolhido;
		\item Compromisso entre vizinhanças pequenas e grandes:
			\begin{itemize}
				\item pequena: iteração barata, mas muitos ótimos locais;
				\item grande: iteração cara, mas menos ótimos locais.
			\end{itemize}
		\item Um ponto que é ótimo local para uma vizinhança pode deixar
			de sê-lo para outra.
	\end{itemize}
\end{frame}

\begin{frame}
	\frametitle{Efeito do tamanho da vizinhança}
	\begin{columns}[T]
	\begin{column}{0.5\textwidth}
	\centering
	{\small Vizinhança pequena $\mathcal{N}_1$}\\[1mm]
	\begin{tikzpicture}[xscale=0.44,yscale=0.6]
		\paisagem
		\fill[red] (1.8,2.6) circle (3pt);
		\draw[red,very thick,|-|] (1.0,-0.5) -- (2.6,-0.5);
		\node[red,below,font=\scriptsize] at (1.8,-0.5) {$\mathcal{N}_1(x)$};
	\end{tikzpicture}\\[1mm]
	{\scriptsize $x$ é ótimo local}
	\end{column}
	\begin{column}{0.5\textwidth}
	\centering
	{\small Vizinhança grande $\mathcal{N}_2$}\\[1mm]
	\begin{tikzpicture}[xscale=0.44,yscale=0.6]
		\paisagem
		\fill[red] (1.8,2.6) circle (3pt);
		\draw[red,very thick,|-|] (0.0,-0.5) -- (6.4,-0.5);
		\node[red,below,font=\scriptsize] at (3.2,-0.5) {$\mathcal{N}_2(x)$};
		\fill[green!50!black] (5.8,3.6) circle (3pt);
	\end{tikzpicture}\\[1mm]
	{\scriptsize $x$ não é ótimo local: o vizinho melhor está em $\mathcal{N}_2$}
	\end{column}
	\end{columns}
\end{frame}

\begin{frame}
	\frametitle{Algoritmo de busca local}
	\begin{enumerate}
		\item Seja $x$ a solução inicial (obtida por uma heurística
			construtiva);
		\item Examinar a vizinhança $\mathcal{N}(x)$;
		\item Se existe $x' \in \mathcal{N}(x)$ com $f(x')$ melhor que
			$f(x)$, fazer $x \leftarrow x'$ e voltar ao passo 2;
		\item Caso contrário, parar: $x$ é um ótimo local em relação a
			$\mathcal{N}$.
	\end{enumerate}
	\vspace{2mm}
	\begin{itemize}
		\item Termina sempre: o custo melhora estritamente a cada passo e
			$S$ é finito;
		\item O resultado depende da solução inicial e da vizinhança.
	\end{itemize}
\end{frame}

\begin{frame}
	\frametitle{Critério de escolha do vizinho}
	\begin{itemize}
		\item \textbf{Primeira aprimorante}: adota o primeiro vizinho
			$x'$ encontrado que melhora o custo de $x$;
			\begin{itemize}
				\item iteração mais barata, caminho de melhoria mais
					longo;
			\end{itemize}
		\item \textbf{Melhor aprimorante}: examina toda a vizinhança
			$\mathcal{N}(x)$ e adota o melhor vizinho $x''$;
			\begin{itemize}
				\item iteração mais cara, passos maiores;
			\end{itemize}
		\item Em ambos os casos, o ponto de parada é um ótimo local em
			relação a $\mathcal{N}$;
		\item Os dois critérios podem terminar em ótimos locais
			diferentes.
	\end{itemize}
\end{frame}

%------------------------------------------------------------
\section{Exploração e explotação}
%------------------------------------------------------------

\begin{frame}
	\frametitle{Exploração e explotação}
	\begin{center}
	\begin{tabular}{p{0.44\textwidth}p{0.44\textwidth}}
		\toprule
		\textbf{Exploração} & \textbf{Explotação} \\
		\midrule
		Visitar regiões distintas do espaço de soluções
		& Examinar com detalhe a região em torno de uma solução \\
		Diversificação
		& Intensificação \\
		Reduz o risco de ignorar a região do ótimo global
		& Atinge o melhor ponto da região atual \\
		Construção aleatória, múltiplas soluções iniciais
		& Busca local \\
		\bottomrule
	\end{tabular}
	\end{center}
\end{frame}

\begin{frame}
	\frametitle{Exploração e explotação na paisagem}
	\begin{columns}[T]
	\begin{column}{0.5\textwidth}
	\begin{itemize}
		\item A busca local é puramente explotatória: sobe até o pico mais
			próximo da solução inicial e para;
		\item Soluções iniciais diferentes podem levar a ótimos locais
			diferentes;
		\item Equilibrar os dois comportamentos é o problema central das
			estratégias mais elaboradas de busca, que ficam fora do
			escopo desta aula.
	\end{itemize}
	\end{column}
	\begin{column}{0.5\textwidth}
	\begin{tikzpicture}[xscale=0.44,yscale=0.6]
		\paisagem
		\fill[orange] (0.5,1.3) circle (3pt);
		\draw[->,orange,thick] (0.55,1.45) -- (1.6,2.45);
		\fill[orange] (4.8,1.9) circle (3pt);
		\draw[->,orange,thick] (4.9,2.0) -- (5.7,3.45);
		\fill[orange] (7.9,1.9) circle (3pt);
		\draw[->,orange,thick] (8.0,2.0) -- (8.9,2.5);
		\node[font=\scriptsize,orange!70!black,align=center] at (4.8,-0.9)
			{três soluções iniciais, três ótimos locais};
	\end{tikzpicture}
	\end{column}
	\end{columns}
\end{frame}

%------------------------------------------------------------
\section{Busca local na mochila}
%------------------------------------------------------------

\begin{frame}
	\frametitle{Mochila: vizinhanças}
	\begin{itemize}
		\item Solução: vetor $x \in \{0,1\}^n$, viável se
			$\sum_i p_i x_i \leq C$;
		\item \textbf{Inversão de bit} (\textit{flip}): troca $x_i$ de
			$0$ para $1$ ou de $1$ para $0$;
			\begin{itemize}
				\item $|\mathcal{N}| = n$;
				\item $0 \to 1$ equivale a \textbf{inserir} um item;
				\item $1 \to 0$ equivale a \textbf{remover} um item;
			\end{itemize}
		\item \textbf{Troca} (\textit{swap}): retira um item que está na
			mochila e insere um que está fora;
			\begin{itemize}
				\item $|\mathcal{N}| = k(n-k)$, com $k$ itens na mochila;
			\end{itemize}
		\item Apenas vizinhos viáveis são considerados.
	\end{itemize}
\end{frame}

\begin{frame}
	\frametitle{Mochila: vizinhança de inversão de bit}
	\begin{center}
	\begin{tikzpicture}
		\node[anchor=east,font=\small] at (0.2,1.0) {solução $x$};
		\vetor{1.0}{1,0,1,1,1,0,0}
		\node[anchor=east,font=\small] at (0.2,0) {inserir item 6};
		\vetor{0}{1,0,1,1,1,1,0}
		\destaque{0}{6}
		\node[anchor=east,font=\small] at (0.2,-0.9) {remover item 3};
		\vetor{-0.9}{1,0,0,1,1,0,0}
		\destaque{-0.9}{3}
	\end{tikzpicture}
	\end{center}
	\begin{itemize}
		\item Inserir aumenta o benefício e o peso: só é aceito se a
			capacidade $C$ for respeitada;
		\item Remover sempre é viável, mas reduz o benefício;
		\item Em um ótimo local para a inversão de bit, nenhum item cabe
			na folga restante.
	\end{itemize}
\end{frame}

\begin{frame}
	\frametitle{Mochila: vizinhança de troca}
	\begin{center}
	\begin{tikzpicture}
		\node[anchor=east,font=\small] at (0.2,0.5) {solução $x$};
		\vetor{0.5}{1,0,1,1,1,0,0}
		\node[anchor=east,font=\small] at (0.2,-0.5) {trocar 4 por 6};
		\vetor{-0.5}{1,0,1,0,1,1,0}
		\destaque{-0.5}{4}
		\destaque{-0.5}{6}
	\end{tikzpicture}
	\end{center}
	\begin{itemize}
		\item Um item sai e outro entra: o peso pode diminuir ou
			aumentar, e o benefício varia pela diferença entre os dois;
		\item Contém movimentos que a inversão de bit não alcança, pois
			combina remoção e inserção em um só passo;
		\item Vizinhança maior que a de inversão de bit: $k(n-k)$
			avaliações por iteração, contra $n$.
	\end{itemize}
\end{frame}

\begin{frame}
	\frametitle{Mochila: instância e solução inicial}
	\begin{columns}[T]
	\begin{column}{0.5\textwidth}
	\begin{center}
	\begin{tabular}{rrrr}
		\toprule
		$i$ & $p_i$ & $b_i$ & $b_i/p_i$ \\
		\midrule
		1 & 4 & 7 & 1,75 \\
		2 & 7 & 6 & 0,86 \\
		3 & 2 & 27 & 13,50 \\
		4 & 7 & 27 & 3,86 \\
		5 & 2 & 19 & 9,50 \\
		6 & 9 & 30 & 3,33 \\
		7 & 9 & 9 & 1,00 \\
		\bottomrule
	\end{tabular}
	\end{center}
	\end{column}
	\begin{column}{0.5\textwidth}
	\begin{itemize}
		\item $n = 7$ e capacidade $C = 18$;
		\item Heurística construtiva gulosa: itens em ordem decrescente de
			$b_i/p_i$, inserindo os que cabem;
		\item Inserção: 3, 5, 4 e 1 (o item 6 não cabe);
		\item $x = (1,0,1,1,1,0,0)$, peso $15$, benefício $80$.
	\end{itemize}
	\end{column}
	\end{columns}
\end{frame}

\begin{frame}
	\frametitle{Mochila: busca local com inversão de bit}
	\begin{itemize}
		\item Solução inicial: $x = (1,0,1,1,1,0,0)$, peso $15$,
			folga $3$, benefício $80$;
		\item Inserir o item 2, 6 ou 7 exigiria peso $22$, $24$ ou
			$24$, acima de $C = 18$: nenhum vizinho viável de inserção;
		\item Remover qualquer item reduz o benefício;
		\item Nenhum vizinho melhora $x$: \textbf{$x$ é ótimo local
			para a inversão de bit}, com benefício $80$.
	\end{itemize}
\end{frame}

\begin{frame}
	\frametitle{Mochila: busca local com troca}
	\begin{columns}[T]
	\begin{column}{0.5\textwidth}
	\begin{center}
	\begin{tabular}{llrr}
		\toprule
		sai & entra & peso & benef. \\
		\midrule
		1 & 2 & 18 & 79 \\
		4 & 2 & 15 & 59 \\
		4 & 6 & 17 & \textbf{83} \\
		4 & 7 & 17 & 62 \\
		\bottomrule
	\end{tabular}
	\end{center}
	{\scriptsize Os outros 8 pares excedem $C = 18$ e são descartados.}
	\end{column}
	\begin{column}{0.5\textwidth}
	\begin{itemize}
		\item Melhor vizinho: trocar o item 4 pelo 6, benefício $83$;
		\item Novo $x = (1,0,1,0,1,1,0)$, peso $17$;
		\item Nova rodada: nenhum vizinho de troca ou de inversão
			melhora $83$.
	\end{itemize}
	\end{column}
	\end{columns}
\end{frame}

\begin{frame}
	\frametitle{Mochila: ótimo local e ótimo global}
	\begin{itemize}
		\item Ótimo local para a inversão de bit: benefício $80$;
		\item Ótimo local para troca (e inversão de bit): benefício $83$;
		\item Ótimo global (enumeração das $2^7$ soluções):
			$x^* = (0,0,1,1,0,1,0)$, peso $18$, benefício $84$;
		\item $x^*$ difere do ótimo local de troca em três posições:
			itens 1, 4 e 5 mudam de estado; nenhuma troca única o alcança.
	\end{itemize}
	\vspace{2mm}
	\begin{center}
	\begin{tikzpicture}
		\node[anchor=east,font=\small] at (0.2,0.5) {inversão de bit};
		\vetor{0.5}{1,0,1,1,1,0,0}
		\node[anchor=west,font=\small] at (5.0,0.5) {$80$};
		\node[anchor=east,font=\small] at (0.2,-0.2) {troca};
		\vetor{-0.2}{1,0,1,0,1,1,0}
		\node[anchor=west,font=\small] at (5.0,-0.2) {$83$};
		\node[anchor=east,font=\small] at (0.2,-0.9) {ótimo global};
		\vetor{-0.9}{0,0,1,1,0,1,0}
		\node[anchor=west,font=\small] at (5.0,-0.9) {$84$};
	\end{tikzpicture}
	\end{center}
\end{frame}

%------------------------------------------------------------
\section{Busca local no TSP: 2-opt}
%------------------------------------------------------------

\begin{frame}
	\frametitle{A vizinhança 2-opt}
	\begin{itemize}
		\item Solução: permutação das cidades (rota);
		\item Movimento: remover duas arestas da rota e reconectar os
			dois caminhos restantes de outro modo, o que inverte o
			trecho entre elas;
		\item $|\mathcal{N}| = n(n-3)/2$ vizinhos por rota;
		\item Remove cruzamentos: toda rota com arestas que se cruzam
			pode ser melhorada por um movimento 2-opt.
	\end{itemize}
	\vspace{1mm}
	\begin{center}
	\begin{tikzpicture}[scale=0.8,
		cid/.style={circle,draw,fill=blue!15,inner sep=1.5pt,font=\scriptsize}]
		% antes
		\begin{scope}
			\foreach \i/\a in {1/90,2/30,3/-30,4/-90,5/-150,6/150}
				\node[cid] (a\i) at (\a:1.2) {\i};
			\draw (a1)--(a2); \draw[red,thick] (a2)--(a5); \draw (a5)--(a4);
			\draw (a4)--(a3); \draw[red,thick] (a3)--(a6); \draw (a6)--(a1);
			\node[font=\small] at (0,-1.9) {antes: 1-2-5-4-3-6-1};
		\end{scope}
		\draw[->,thick] (1.9,0) -- (3.0,0);
		% depois
		\begin{scope}[xshift=5cm]
			\foreach \i/\a in {1/90,2/30,3/-30,4/-90,5/-150,6/150}
				\node[cid] (b\i) at (\a:1.2) {\i};
			\draw (b1)--(b2); \draw[green!50!black,thick] (b2)--(b3);
			\draw (b3)--(b4); \draw (b4)--(b5);
			\draw[green!50!black,thick] (b5)--(b6); \draw (b6)--(b1);
			\node[font=\small] at (0,-1.9) {depois: 1-2-3-4-5-6-1};
		\end{scope}
	\end{tikzpicture}
	\end{center}
\end{frame}

\begin{frame}
	\frametitle{Variação de custo de um movimento 2-opt}
	\begin{itemize}
		\item Rota com arestas consecutivas $(a,b)$ e $(c,d)$, nesta ordem;
		\item O movimento substitui $(a,b)$ e $(c,d)$ por $(a,c)$ e
			$(b,d)$ e inverte o trecho de $b$ a $c$;
		\item Variação de custo, calculada em tempo constante:
			\[
				\Delta = d_{ac} + d_{bd} - d_{ab} - d_{cd}
			\]
		\item Se $\Delta < 0$, o movimento melhora a rota;
		\item O custo total não precisa ser recalculado: basta somar
			$\Delta$ ao custo atual.
	\end{itemize}
\end{frame}

\begin{frame}[fragile]
	\frametitle{Busca local 2-opt}
	\begin{lstlisting}[language=Python]
melhorou = True
while melhorou:
    melhorou = False
    for i in range(0, n - 2):
        for j in range(i + 2, n):
            if i == 0 and j == n - 1:
                continue
            a, b = rota[i], rota[i + 1]
            c, e = rota[j], rota[(j + 1) % n]
            delta = dist[a][c] + dist[b][e] - dist[a][b] - dist[c][e]
            if delta < 0:
                rota[i + 1 : j + 1] = reversed(rota[i + 1 : j + 1])
                custo += delta
                melhorou = True
	\end{lstlisting}
	\begin{itemize}
		\item Variante de primeira aprimorante: aplica o movimento assim
			que $\Delta < 0$;
		\item Variante de melhor aprimorante: percorre todos os pares e
			aplica só o de menor $\Delta$;
		\item A rota ao final é um ótimo local em relação ao 2-opt.
	\end{itemize}
\end{frame}

\begin{frame}
	\frametitle{Fluxo completo: construção e busca local}
	\begin{itemize}
		\item Heurística construtiva (vizinho mais próximo, multi-start
			ou aleatória) fornece a solução inicial;
		\item A busca local 2-opt leva a solução a um ótimo local;
		\item A qualidade final depende da solução inicial e da
			vizinhança;
		\item O ótimo local não é, em geral, o ótimo global: o GAP em
			relação ao ótimo conhecido costuma diminuir, mas não se
			anula.
	\end{itemize}
\end{frame}

%------------------------------------------------------------
\section{A tarefa (avaliativa)}
%------------------------------------------------------------

\begin{frame}
	\frametitle{Caráter avaliativo}
	\begin{itemize}
		\item Esta tarefa vale \textbf{25\% da nota da Unidade 1};
		\item Pode ser realizada em equipes de até 3 pessoas;
		\item Prazo: sexta-feira, \textbf{09/10}, às 8h50;
		\item Prossegue a implementação do TSP: as quatro instâncias, os
			ótimos conhecidos e as três heurísticas construtivas da
			tarefa anterior são reutilizados.
	\end{itemize}
\end{frame}

\begin{frame}
	\frametitle{O que é pedido}
	\begin{itemize}
		\item Implementar a vizinhança 2-opt e a busca local
			correspondente, com variação de custo calculada em tempo
			constante;
		\item Aplicar a busca local à solução de cada heurística
			construtiva: vizinho mais próximo (cidade 1), vizinho mais
			próximo multi-start e construção aleatória (30 execuções);
		\item Montar uma tabela única com, por instância e heurística, o
			GAP da solução inicial e o GAP da solução após o 2-opt, em
			relação ao ótimo conhecido;
		\item Descrever o ambiente computacional e analisar os
			resultados.
	\end{itemize}
\end{frame}

%------------------------------------------------------------
\section{Referências}
%------------------------------------------------------------

\begin{frame}
	\frametitle{Referências}
	\begin{itemize}
		\item Goldbarg, M. C.; Luna, H. P. L. \textit{Otimização
			Combinatória e Programação Linear}. 2005.
		\item Hillier, F. S.; Lieberman, G. J. \textit{Introduction to
			Operations Research}. 2010.
		\item Croes, G. A. A method for solving traveling-salesman
			problems. \textit{Operations Research}, 6(6), 1958.
	\end{itemize}
\end{frame}

\end{document}
